\chapter{Design Decision Log}
\label{app:design-decisions}

This appendix catalogues the non-obvious architectural choices made
during the modernisation of the Benton--Kennedy domain-theory library for
Rocq~9.0.  Each entry records the decision, the alternatives considered,
and the rationale that drove the final choice.  Entries are numbered
sequentially in the order they were taken; cross-references to the
relevant source files and thesis sections are given where appropriate.

% ──────────────────────────────────────────────────────────────────
\section*{DD-001: No separate \texttt{Pointed.v}}

\textbf{Decision.}
\texttt{HasBottom}, \texttt{IsPointed}, and \texttt{PointedCPO} live in
\texttt{CPO.v}; \texttt{strict\_fun} lives in \texttt{Morphisms.v}.
No \texttt{Pointed.v} file is created.

\textbf{Rationale.}
An earlier iteration had a one-line re-export shim that bounced users
from \texttt{Pointed.v} back to \texttt{CPO.v}.  Following the
MathComp/HB convention of organising files around dependency structure
rather than an idealised concept map, pointedness belongs in
\texttt{CPO.v} (hard dependency on \texttt{CPO~T}) and strictness in
\texttt{Morphisms.v} (hard dependency on \texttt{cont\_fun}).

% ──────────────────────────────────────────────────────────────────
\section*{DD-002: Separate continuity axioms in \texttt{IsCPOEnriched}}

\textbf{Decision.}
\texttt{IsCPOEnriched} axiomatises composition via two separate
Scott-continuity conditions (\texttt{comp\_cont\_l} and
\texttt{comp\_cont\_r}) rather than a single joint-continuity condition.

\textbf{Rationale.}
Joint continuity is equivalent to separate continuity in each variable
(Abramsky \& Jung, Lemma~3.2.6).  The separate form avoids importing
\texttt{Products.v} at the point where \texttt{Enriched.v} is declared,
breaking the dependency cycle
\texttt{Enriched.v $\to$ Products.v $\to$ CPOTheory.v $\to$ Enriched.v}.
This is the same strategy used by
Kornell--Lindenhovius--Mislove~\cite{KLM2024}~\S3.3.

% ──────────────────────────────────────────────────────────────────
\section*{DD-003: Monotonicity fields precede continuity fields}

\textbf{Decision.}
In every HB mixin that requires both monotonicity and Scott-continuity,
the monotonicity field is declared \emph{before} the continuity field.

\textbf{Rationale.}
The \texttt{continuous} predicate takes a \texttt{mono\_fun} as its
argument; \texttt{Build\_mono\_fun} requires a proof of monotonicity.
If the continuity field appeared first, the monotonicity proof would not
yet be in scope.  This is a Rocq/HB implementation detail, not a
mathematical preference.

% ──────────────────────────────────────────────────────────────────
\section*{DD-004: Leibniz equality for categorical laws}

\textbf{Decision.}
The categorical laws in \texttt{IsCPOEnriched} and
\texttt{IsLocallyContinuous} use Leibniz equality \texttt{=} rather than
setoid equality \texttt{$\approx$}.

\textbf{Rationale.}
The hom-CPOs are HB-packed \texttt{CPO.type} values whose element
equality is Leibniz by default.  Using \texttt{=} keeps the laws
definitionally true and avoids setoid-rewriting overhead at the structure
level.  If a concrete instance uses a coarser extensional equality, the
laws can be restated in the relevant theory file.

% ──────────────────────────────────────────────────────────────────
\section*{DD-005: \texttt{HasMixedEndo} records only data}

\textbf{Decision.}
The mixin \texttt{HasMixedEndo} in \texttt{Enriched.v} contains only the
object and morphism actions of a mixed-variance bifunctor.  The
functoriality and local-continuity axioms are deferred to
\texttt{DomainEquations.v}.

\textbf{Rationale.}
The full axioms require reasoning about directed sets of morphism pairs,
which in turn requires the product CPO construction.  Declaring them in
\texttt{Enriched.v} would create the same forward-dependency cycle
described in DD-002.  Splitting data from axioms keeps the structures
file dependency-minimal.

% ──────────────────────────────────────────────────────────────────
\section*{DD-006: Classical lift via \texttt{option~D}}

\textbf{Decision.}
The main library lift (\texttt{Lift.v}) uses the flat carrier
\texttt{option~D} with a classically computed sup.  A coinductive
alternative is developed in the supplementary file \texttt{LiftMonad.v}.

\textbf{Rationale.}
The coinductive \texttt{delay~D} type suffers an \emph{antisymmetry
obstruction}: \texttt{now~d} and \texttt{later~(now~d)} are
propositionally distinct yet semantically equivalent, so any CPO order
compatible with convergence would force
\texttt{now~d $=$ later~(now~d)}---a contradiction.  Making
\texttt{delay~D} into a \texttt{CPO.type} (which uses Leibniz equality,
DD-004) therefore requires quotienting by weak bisimulation, which in
turn requires either a setoid-based CPO framework or quotient types not
standard in Rocq.  The \texttt{option~D} approach avoids the quotient
entirely; its only non-constructive step is \texttt{ClassicalEpsilon},
localised to \texttt{lift\_sup}.

% ──────────────────────────────────────────────────────────────────
\section*{DD-007: \texttt{LiftMonad.v} as a self-contained supplementary file}

\textbf{Decision.}
The coinductive delay monad is developed in a separate, zero-import file
\texttt{LiftMonad.v} that is not used by any other file in the library.

\textbf{Rationale.}
The file serves two thesis purposes: (i) concrete comparison of the
\texttt{option~D} and \texttt{delay~D} monad structures side by side;
(ii) a precise Rocq proof of \texttt{delay\_CPO\_requires\_quotient},
which establishes that \texttt{delay~D} cannot be a \texttt{CPO.type}
without quotient types.  Keeping the file self-contained emphasises the
logical independence of the coinductive structure from the domain-theory
hierarchy.

% ──────────────────────────────────────────────────────────────────
\section*{DD-008: HB instance ordering for \texttt{PointedCPO}}

\textbf{Decision.}
For any type~$T$, the instance registrations \texttt{lift\_IsCPO},
\texttt{lift\_HasBottom}, \texttt{lift\_IsPointed} must appear in exactly
this order.

\textbf{Rationale.}
\texttt{IsPointed} requires both \texttt{CPO~T} and
\texttt{HasBottom~T}.  Violating the order causes a silent HB
elaboration failure (no error; the instance simply does not register).
The general rule: \texttt{HasLe $\to$ IsPreorder $\to$ IsPartialOrder
$\to$ HasSup $\to$ IsCPO $\to$ HasBottom $\to$ IsPointed}.

% ──────────────────────────────────────────────────────────────────
\section*{DD-009: Intrinsic typing and ANF for PCF\textsubscript{v}}

\textbf{Decision.}
PCF syntax uses intrinsically typed terms in Administrative Normal Form,
with separate mutually inductive \texttt{Value} and \texttt{Exp} types
indexed by \texttt{Env} and \texttt{Ty}.  Variables are typed de~Bruijn
indices; renamings bootstrap substitutions
(\cref{sec:meth-pcf-syntax}).

\textbf{Rationale.}
Benton--Kennedy~\S3 describe two approaches: extrinsic typing (raw terms
plus a separate typing judgment) and intrinsic typing (well-typedness by
construction).  The extrinsic approach required lengthy proofs that any
two typing derivations are equal, plus side-conditions on every
definition.  The intrinsic approach eliminates both problems.  ANF and
the renaming bootstrap follow Benton--Kennedy exactly; a complete naming
comparison appears in \cref{sec:disc-bkv}.

% ──────────────────────────────────────────────────────────────────
\section*{DD-010: Big-step CBV evaluation}

\textbf{Decision.}
The operational semantics uses a big-step evaluation relation
\texttt{Eval : CExp\,$\tau$ $\to$ CValue\,$\tau$ $\to$ Prop}
(notation $e \Downarrow v$), not a small-step reduction relation.

\textbf{Rationale.}
Big-step evaluation matches the denotational semantics for CBV: the
soundness theorem has the form
$e \Downarrow v \;\Rightarrow\; \llbracket e \rrbracket = \eta(\llbracket v \rrbracket)$.
With intrinsically typed closed terms there are no stuck states, so no
progress or preservation theorems are needed.

% ──────────────────────────────────────────────────────────────────
\section*{DD-011: Point-free mutual fixpoint for denotation}

\textbf{Decision.}
The denotation functions \texttt{sem\_val} and \texttt{sem\_exp} are
defined as a single mutual structural \texttt{Fixpoint}, with each case
expressed as a point-free composition of library combinators
(\cref{sec:meth-pcf-denotational}).

\textbf{Rationale.}
Two approaches were evaluated: (i)~a three-pass approach (raw functions,
then monotonicity, then continuity); (ii)~a single-pass point-free
style where each case is a pipeline of combinators already packaged as
\texttt{cont\_fun}.  The single-pass approach was chosen because the
combinator library already provides all needed packaging and the
\texttt{FIX} case benefits from \texttt{FIXP} being a single continuous
operator.  The downside is that \texttt{simpl} produces unreadable
goals; this is mitigated by using \texttt{change} rather than
\texttt{simpl} and making \texttt{sem\_ren} \texttt{Opaque} during the
renaming/substitution proofs.

% ──────────────────────────────────────────────────────────────────
\section*{DD-012: Renaming route for the substitution lemma}

\textbf{Decision.}
The substitution lemma is proved via an intermediate renaming route:
first the renaming lemma by mutual induction, then the weakening
corollary, then the extension lemma, and finally the substitution
lemma by mutual induction (\cref{sec:meth-pcf-denotational}).

\textbf{Rationale.}
The direct approach to the substitution extension lemma requires knowing
that syntactic weakening (\texttt{wkVal}) corresponds to dropping the
first environment component.  Since \texttt{wkVal\,v $=$ renVal\,ren\_wk\,v},
this requires the renaming lemma first.  The dependency chain mirrors
Benton--Kennedy's syntactic bootstrapping strategy.

% ──────────────────────────────────────────────────────────────────
\section*{DD-013: \texttt{var\_case} combinator}

\textbf{Decision.}
The renaming and substitution extension functions use a
\texttt{var\_case} combinator that eliminates the head variable via a
dependent match with a function-returning return type, rather than
pattern matching that introduces \texttt{JMeq\_eq}
(\cref{sec:meth-pcf-syntax}).

\textbf{Rationale.}
A na\"{\i}ve dependent match on \texttt{Var\,($\tau$\,::\,$\Gamma$)\,$\sigma$}
introduces the opaque axiom \texttt{JMeq\_eq} from
\texttt{Program.Equality}, which blocks kernel reduction and causes
\texttt{reflexivity} to fail on computation rules.  The
\texttt{var\_case} combinator avoids this by using a return type that
takes the head and tail handlers as parameters \emph{after} the match,
so that \texttt{var\_case\,z\,s\,ZVAR $\equiv$ z} and
\texttt{var\_case\,z\,s\,(SVAR\,x') $\equiv$ s\,\_\,x'} hold by pure
$\iota$+$\beta$ reduction.

% ──────────────────────────────────────────────────────────────────
\section*{DD-014: Soundness by structural induction on \texttt{Eval}}

\textbf{Decision.}
Soundness is proved by structural induction on the big-step evaluation
derivation, using computation rules and substitution lemmas directly.

\textbf{Rationale.}
In the CBV big-step setting, each evaluation rule corresponds to a
computation rule plus induction hypotheses.  The two non-trivial cases
(\texttt{LET} and \texttt{APP}) require only \texttt{sem\_single\_subst}
and \texttt{sem\_double\_subst}.  No auxiliary lemmas outside
\texttt{PCF\_Soundness.v} are needed.

% ──────────────────────────────────────────────────────────────────
\section*{DD-015: Adequacy via logical relation}

\textbf{Decision.}
Computational adequacy is proved via a type-indexed logical relation
following Benton--Kennedy--Varming~\S3.2, with the
\texttt{FIX} case appealing to \texttt{fixp\_ind}.

\textbf{Rationale.}
The relation separates into \texttt{rel\_val} (values) and
\texttt{rel\_exp} (expressions, the ``lift'' of \texttt{rel\_val}).  The
arrow case uses \texttt{doubleSubstExp} to mirror the operational
\texttt{e\_App} rule directly.  Admissibility of \texttt{rel\_exp} is
established by \texttt{rel\_val\_admissible} (induction on~$\tau$) and
chain-stabilisation properties of the lift CPO.  The fundamental lemma
is proved by mutual induction on \texttt{Value}/\texttt{Exp} via the
\texttt{Combined Scheme val\_exp\_ind}.

% ──────────────────────────────────────────────────────────────────
\section*{DD-016: HB coercion workarounds in \texttt{EnrichedTheory.v}}

\textbf{Decision.}
\texttt{EnrichedTheory.v} uses a systematic set of HB coercion
workarounds---\texttt{etransitivity}, explicit \texttt{@}-qualified
projections, \texttt{unshelve econstructor}, \texttt{Arguments\,\{C\}}
directives, and \texttt{cbn\,[mf\_fun]} normalisation---to compile from
source rather than from \texttt{.vo} files.

\textbf{Rationale.}
When building from source, HB canonical-structure resolution fails in
several contexts that work from \texttt{.vo} files in the IDE.  The core
problem is that coercion paths between \texttt{CPOEnrichedCat.sort} and
\texttt{CPO.sort} resolve ambiguously.  Five categories of workarounds
were identified (\cref{sec:meth-enriched}).  Additionally, joint
continuity of composition is proved in a two-stage approach: a
product-free stage~A that avoids coercion-path conflicts, and a thin
packaging stage~B that confines the conflicts.

% ──────────────────────────────────────────────────────────────────
\section*{DD-017: \texttt{IsMixedLocallyContinuous} axiom design}

\textbf{Decision.}
The \texttt{IsMixedLocallyContinuous} mixin axiomatises the
mixed-variance bifunctor via six separate fields---identity,
composition, separate monotonicity, and separate continuity---rather
than joint conditions.

\textbf{Rationale.}
This mirrors DD-002: separate axioms avoid importing
\texttt{Products.v} at the mixin declaration point.  Joint monotonicity
and continuity are derived in a \texttt{MFDerived} section.
\texttt{DomainEquations.v} also contains the axiom
\texttt{bilimit\_exists}, which asserts the existence of the bilimit of
the approximation sequence; its proof requires an $\omega$-product CPO
not yet in the library.  All consequences of the bilimit (\texttt{ROLL},
\texttt{UNROLL}, deflation chains) are fully proved from the record
fields.

% ──────────────────────────────────────────────────────────────────
\section*{DD-018: \texttt{nat\_trans} as a plain record}

\textbf{Decision.}
Natural transformations are defined as a plain \texttt{Record nat\_trans}
parameterised over two \texttt{lc\_functor} values, not over HB
\texttt{LocallyContinuousFunctor} instances.

\textbf{Rationale.}
HB instances for functors on \texttt{CPOEnrichedCat.type} trigger
universe inconsistencies when natural transformations quantify over all
objects.  Plain records avoid HB's universe machinery and compose
directly (identity, vertical, horizontal composition produce new
\texttt{nat\_trans} values without instance registration).

% ──────────────────────────────────────────────────────────────────
\section*{DD-019: Yoneda lemma in \texttt{instances/}}

\textbf{Decision.}
The representable functor and the enriched Yoneda lemma live in
\texttt{instances/Yoneda.v} rather than in \texttt{NatTrans.v}.

\textbf{Rationale.}
The representable functor requires the concrete CPO-enriched category
instance registered in \texttt{instances/Function.v}.  Placing Yoneda in
\texttt{NatTrans.v} (a theory file) would create a reverse dependency
from \texttt{theory/} to \texttt{instances/}, violating the library's
build-order constraint.

% ──────────────────────────────────────────────────────────────────
\section*{DD-020: \texttt{change} tactic for \texttt{FunLift.v}}

\textbf{Decision.}
The axiom proofs in \texttt{FunLift.v} use the \texttt{change} tactic
to normalise coercion-wrapped goals, rather than \texttt{rewrite} or
\texttt{unfold}.

\textbf{Rationale.}
After case-splitting on \texttt{option} values, goals are wrapped in
HB-generated coercion chains that prevent \texttt{rewrite} from matching
the LHS.  \texttt{change} bypasses coercions entirely by checking
definitional equality directly.  Both \texttt{unfold + rewrite} and
\texttt{subst + unfold} strategies were tried and fail due to immediate
$\iota$-reduction that eliminates the head symbol before
\texttt{rewrite} can match.

% ──────────────────────────────────────────────────────────────────
\section*{DD-022: Atoms-only quantum representation}

\textbf{Decision.}
Quantum sets are represented as plain Rocq types (atom indices) with
$Q$-valued relations, where $Q$ is axiomatised as an involutive quantale
via HB.  Quantum posets are plain \texttt{Record qposet} values
parameterised by~$Q$, not HB structures.

\textbf{Rationale.}
KLM's quantum CPO theory separates into two layers: (i) abstract order
theory (convergence, completeness, Scott continuity, enrichment), which
needs only a quantale~$Q$ and $Q$-valued relations; (ii) concrete
operator-space constructions (tensor products, monoidal closure), which
require Hilbert spaces and completely bounded maps.  Layer~(ii) does not
exist in any Rocq library.  The atoms-only approach covers Layer~(i)
completely, matches Rocq's strengths in abstract algebra, and is
justified by the existence of $B(H)$ as a concrete involutive quantale.
Quantum posets are plain records (not HB structures) because HB does not
handle structures parameterised by other structures well, following the
same pattern as \texttt{chain}, \texttt{mono\_fun}, and
\texttt{cont\_fun}.
